%TC:ignore - ignoriert die folgnden Wörter beim Wordcount in Overleaf 
%-----------------------------------
% Define document and include general packages
%-----------------------------------
% Tabellen- und Abbildungsverzeichnis stehen normalerweise nicht im
% Inhaltsverzeichnis. Gleiches gilt für das Abkürzungsverzeichnis (siehe unten).
% Manche Dozenten bemängeln das. Die Optionen 'listof=totoc,bibliography=totoc'
% geben das Tabellen- und Abbildungsverzeichnis im Inhaltsverzeichnis (toc=Table
% of Content) aus.
% Da es aber verschiedene Regelungen je nach Dozent geben kann, werden hier
% beide Varianten dargestellt.
\documentclass[12pt,oneside,titlepage,listof=totoc,bibliography=totoc,listof=entryprefix]{scrartcl}
%\documentclass[12pt,oneside,titlepage]{scrartcl}

%-----------------------------------
% Dokumentensprache
%-----------------------------------
%\def\FOMEN{}% Auskommentieren um die Dokumentensprache auf englisch zu ändern
\newif\ifde
\newif\ifen
\newif\ifshowListOfAppendices
\showListOfAppendicestrue
% \showListOfAppendicesfalse

%-----------------------------------
% Meta informationen
%-----------------------------------
%-----------------------------------
% Meta Informationen zur Arbeit
%-----------------------------------

% Autor
\newcommand{\myAutor}{Autor:in 1}
\newcommand{\myAutorzwei}{Autor:in 2}
\newcommand{\myAutordrei}{Autor:in 3}
\newcommand{\myAutorvier}{Autor:in 4}

% Adresse
\newcommand{\myAdresse}{Heidestra\ss e 17 \\ \> \> \> 51147 Köln}

% Titel der Arbeit
\newcommand{\myTitel}{\LaTeX}

% Betreuer
\newcommand{\myBetreuer}{Prof. Dr. Peter Lustig}

% Lehrveranstaltung
\newcommand{\myModul}{Modul Nr. 1}

% Matrikelnummer
\newcommand{\myMatrikelNr}{123456}
\newcommand{\myMatrikelNrzwei}{234567}
\newcommand{\myMatrikelNrdrei}{234567}
\newcommand{\myMatrikelNrvier}{234567}

% Ort
\newcommand{\myOrt}{Berlin}
\newcommand{\myOrtzwei}{München}
\newcommand{\myOrtdrei}{Berlin}
\newcommand{\myOrtvier}{München}

% Datum der Abgabe
\newcommand{\myAbgabeDatum}{\today}

% Semesterzahl
\newcommand{\mySemesterZahl}{4}

% Name der Hochschule
\newcommand{\myHochschulName}{FOM Hochschule für Oekonomie \& Management}

% Standort der Hochschule
\newcommand{\myHochschulStandort}{Dortmund}

% Studiengang
\newcommand{\myStudiengang}{mein Studiengang}

% Art der Arbeit
\newcommand{\myThesisArt}{Hausarbeit}

% Zu erlangender akademische Grad
\newcommand{\myAkademischerGrad}{Bachelor of Science (B.Sc.)}

% Firma
\newcommand{\myFirma}{Mustermann GmbH}

% Definition der Sprache
\ifdefined\FOMEN
%Englisch
\entrue
\usepackage[english]{babel}
\else
%Deutsch
\detrue
\usepackage[ngerman]{babel}
\fi



\newcommand{\langde}[1]{%
   \ifde\selectlanguage{ngerman}#1\fi}
\newcommand{\langen}[1]{%
   \ifen\selectlanguage{english}#1\fi}
\usepackage[utf8]{luainputenc}
\langde{\usepackage[babel,german=quotes]{csquotes}}
\langen{\usepackage[babel,english=british]{csquotes}}
\ifde\usepackage[ngerman]{babel}\fi
\usepackage{fancyhdr}
\usepackage[german, useregional=numeric]{datetime2} %einfachere Datumsangaben
\usepackage{fancybox}
\usepackage[a4paper, left=4cm, right=2cm, top=3.9cm, bottom=2cm, marginparwidth=3.5cm, marginparsep=0.3cm, headheight=14.5pt, headsep=1.5cm]{geometry}
\usepackage{graphicx}
\usepackage{colortbl}
\renewcommand\familydefault{\sfdefault}
\usepackage{ragged2e}
\usepackage[capposition=top]{floatrow}
\floatsetup{capposition=top, justification=raggedright}
\usepackage{array}
\usepackage{float}      %Positionierung von Abb. und Tabellen mit [H] erzwingen
\usepackage{footnote}
% Darstellung der Beschriftung von Tabellen und Abbildungen (Leitfaden S. 44)
% singlelinecheck=false: macht die Caption linksbündig (statt zentriert)
% labelfont auf fett: (Tabelle x.y:, Abbildung: x.y)
% font auf fett: eigentliche Bezeichnung der Abbildung oder Tabelle
% Fettschrift laut Leitfaden 2018 S. 45
\usepackage{caption}
\captionsetup{format=plain, singlelinecheck=false, justification=RaggedRight, labelfont=bf, font=bf, aboveskip=12pt, belowskip=6pt}
\usepackage{enumitem}
\usepackage{amssymb}
\usepackage{mathptmx}
\usepackage{minted} %Kann für schöneres Syntax Highlighting genutzt werden. ACHTUNG: Python muss installiert sein.
\usepackage{amsmath}
\usepackage{lmodern}
\usepackage[T1]{fontenc} %% Schriftartoptionen:
%\usepackage[scaled=0.91]{helvet} % Behebt, zusammen mit Package courier, pixelige Überschriften. Ist, zusammen mit mathptx, dem times-Package vorzuziehen. Details: https://latex-kurs.de/fragen/schriftarten/Times_New_Roman.html
%\usepackage{courier}
%\usepackage{inter}
%\renewcommand{\familydefault}{\sfdefault}
\usepackage{fontspec}
\setmainfont[Scale=0.9]{TeX Gyre Heros} % Helvetica Klon 
\setmonofont{Courier New}
\usepackage[mathrm=sym]{unicode-math}
%\setmathfont{XITS Math}
%\setmathfont{Latin Modern Math}
\setmathfont{TeX Gyre Pagella Math} % Schola, Termes, Bonum

\usepackage[table]{xcolor}
\usepackage{marvosym}% Verwendung von Symbolen, z.B. perfektes Eurozeichen
\usepackage{wrapfig} % Bilder im Text mit Textumbruch

\usepackage{changepage} % Für die korrekte Einrückung von 1cm bei Formeln

%-----------------------------------
% Formelverzeichnis
%-----------------------------------
\DeclareNewTOC[%
tocentryindent   = 1.5em,
tocentrynumwidth = 5em,
name=Formel,
type=xequation,
nonfloat
]{loe}

\AtBeginDocument{%
\newcaptionname{ngerman}\xequationname{Formel}%
\newcaptionname{ngerman}\listxequationname{Formel- und Symbolverzeichnis}%
}


%-----------------------------------
% Überschriften auf 12pt setzen
% Absatz – 12pt davor, 6pt danach
%-----------------------------------
\addtokomafont{section}{\fontsize{12}{21.6}\selectfont}
\addtokomafont{subsection}{\fontsize{12}{21.6}\selectfont}
\addtokomafont{subsubsection}{\fontsize{12}{21.6}\selectfont}
\addtokomafont{paragraph}{\fontsize{12}{21.6}\selectfont}
\addtokomafont{subparagraph}{\fontsize{12}{21.6}\selectfont}

% Abstände vor (beforeskip) und nach (afterskip) Überschriften
% Negativer beforeskip-Wert verhindert Einrückung der ersten Zeile nach der Überschrift
\RedeclareSectionCommand[beforeskip=12pt, afterskip=6pt]{section}
\RedeclareSectionCommand[beforeskip=12pt, afterskip=6pt]{subsection}
\RedeclareSectionCommand[beforeskip=12pt, afterskip=6pt]{subsubsection}
\RedeclareSectionCommand[beforeskip=12pt, afterskip=6pt]{paragraph}
\RedeclareSectionCommand[beforeskip=12pt, afterskip=6pt]{subparagraph}


%% Anpassen der Überschriften im Inhaltsverzeichnis
% Abstände zwischen den Überschriften können hier für jede Ebene separat bestimmt werden
\RedeclareSectionCommand[tocbeforeskip=0pt]{section}
\RedeclareSectionCommand[tocbeforeskip=0pt]{subsection}
\RedeclareSectionCommand[tocbeforeskip=0pt]{subsubsection}
\RedeclareSectionCommand[tocbeforeskip=0pt]{paragraph}
\RedeclareSectionCommand[tocbeforeskip=0pt]{subparagraph}

%section Überschriften werden hier in normaler Schrift dargestellt und es werden linefill Punkte hinzugefügt
\DeclareTOCStyleEntry[
  entryformat=\normalfont,
  pagenumberformat=\normalfont,
  linefill=\TOCLineLeaderFill
]{tocline}{section}

\DeclareTOCStyleEntry[
  entryformat=\normalfont,
  pagenumberformat=\normalfont
]{tocline}{subsection}

\DeclareTOCStyleEntry[
  entryformat=\normalfont,
  pagenumberformat=\normalfont
]{tocline}{subsubsection}

\newlength{\originalintextsep}
\setlength{\originalintextsep}{\intextsep} % \intextsep kann mit \setlength{\intextsep}{0pt} auf 0pt gesetzt werden z.B. vor wrapfigure % Margen über und unter der wrapfigure entfernen
% reset mit \setlength{\intextsep}{\originalintextsep}

% Mehrere Fussnoten nacheinander mit Komma separiert
\usepackage[hang,multiple]{footmisc}
\setlength{\footnotemargin}{1em}
\setlength{\skip\footins}{1.5em plus 0.3em minus 0.5em}

% todo Aufgaben als Kommentare verfassen für verschiedene Editoren
\usepackage[textsize=tiny, colorinlistoftodos, tickmarkheight=5pt]{todonotes}
\setlength{\marginparwidth}{3.5cm}
\reversemarginpar
% Verhindert, dass nur eine Zeile auf der nächsten Seite steht
%\setlength{\marginparwidth}{2cm}
\usepackage[all]{nowidow}

%-----------------------------------
% Farbdefinitionen
%-----------------------------------
\definecolor{darkblack}{rgb}{0,0,0}
\definecolor{dunkelgrau}{rgb}{0.8,0.8,0.8}
\definecolor{hellgrau}{rgb}{0.0,0.7,0.99}
\definecolor{mauve}{rgb}{0.58,0,0.82}
\definecolor{dkgreen}{rgb}{0,0.6,0}

%-----------------------------------
% Pakete für Tabellen
%-----------------------------------
\usepackage{epstopdf}
\usepackage{nicefrac} % Brüche
\usepackage{multirow}
\usepackage{rotating} % vertikal schreiben
\usepackage{mdwlist}
\usepackage{tabularx}% für Breitenangabe

\setlength{\originalintextsep}{\intextsep}

%-----------------------------------
% sauber formatierter Quelltext
%-----------------------------------
\usepackage{listings}
% JavaScript als Sprache definieren:
\lstdefinelanguage{JavaScript}{
	keywords={break, super, case, extends, switch, catch, finally, for, const, function, try, continue, if, typeof, debugger, var, default, in, void, delete, instanceof, while, do, new, with, else, return, yield, enum, let, await},
	keywordstyle=\color{blue}\bfseries,
	ndkeywords={class, export, boolean, throw, implements, import, this, interface, package, private, protected, public, static},
	ndkeywordstyle=\color{darkgray}\bfseries,
	identifierstyle=\color{black},
	sensitive=false,
	comment=[l]{//},
	morecomment=[s]{/*}{*/},
	commentstyle=\color{purple}\ttfamily,
	stringstyle=\color{red}\ttfamily,
	morestring=[b]',
	morestring=[b]"
}

\lstset{
	%language=JavaScript,
	numbers=left,
	numberstyle=\tiny,
	numbersep=5pt,
	breaklines=true,
	showstringspaces=false,
	frame=l,
	xleftmargin=5pt,
	xrightmargin=5pt,
	basicstyle=\ttfamily\scriptsize,
	stepnumber=1,
	keywordstyle=\color{blue},          % keyword style
  	commentstyle=\color{dkgreen},       % comment style
  	stringstyle=\color{mauve}         % string literal style
}

%-----------------------------------
%Literaturverzeichnis Einstellungen
%-----------------------------------

% Biblatex

\usepackage{url}
\urlstyle{same}
\newcommand{\citationstyle}{fom_2018} % Mögliche Werte: ieee, fom_2018, fom_alt
% Laden des entsprechenden Zitationsstils basierend auf der Variable \citationstyle
\ifthenelse{\equal{\citationstyle}{ieee}}{
    % Einstellungen für IEEE-Zitationsstil
	\usepackage[
		backend=biber,
		style=ieee,
		maxcitenames=3,	% mindestens 3 Namen ausgeben bevor et. al. kommt
		maxbibnames=999,
		date=iso,
		seconds=true, %werden nicht verwendet, so werden aber Warnungen unterdrückt.
		urldate=iso,
		dashed=false,
		autocite=inline,
		useprefix=true, % 'von' im Namen beachten (beim Anzeigen)
		mincrossrefs = 1
	]{biblatex}%iso dateformat für YYYY-MM-DD
    % et al. anstatt u. a. bei mehr als drei Autoren.
    \DefineBibliographyStrings{ngerman}{
        andothers = {{et\,al\adddot}},
    }
    \DefineBibliographyStrings{english}{
        andothers = {{et\,al\adddot}},
    }
}{
    \ifthenelse{\equal{\citationstyle}{fom_2018}}{
        % Einstellungen für Neuer Leitfaden (2018)
		\usepackage[
			backend=biber,
			style=ext-authoryear-ibid, % Auskommentieren und nächste Zeile einkommentieren, falls "Ebd." (ebenda) nicht für sich-wiederholende Fussnoten genutzt werden soll.
			%style=ext-authoryear,
			maxcitenames=3,	% mindestens 3 Namen ausgeben bevor et. al. kommt
			maxbibnames=999,
			mergedate=false,
			date=iso,
			seconds=true, %werden nicht verwendet, so werden aber Warnungen unterdrückt.
			urldate=iso,
			innamebeforetitle,
			dashed=false,
			autocite=footnote,
			doi=false,
			useprefix=true, % 'von' im Namen beachten (beim Anzeigen)
			mincrossrefs = 1
		]{biblatex}%iso dateformat für YYYY-MM-DD
        %weitere Anpassungen für BibLaTex
        \input{skripte/modsBiblatexZotero}
    }{
        \ifthenelse{\equal{\citationstyle}{fom_alt}}{
            % Einstellungen für Alter Leitfaden
            \usepackage[
                backend=biber,
                style=numeric,
                citestyle=authoryear,
                url=false,
                isbn=false,
                notetype=footonly,
                hyperref=false,
                sortlocale=de
            ]{biblatex}
            %weitere Anpassungen für BibLaTex
            \input{skripte/modsBiblatex}
        }{
            % Fehlerbehandlung für ungültige Werte
            \PackageError{thesis_main}{Ungueltiger Wert fuer \noexpand\citationstyle: \citationstyle}{Bitte setzen Sie \noexpand\citationstyle auf ieee, fom_2018 oder fom_alt.}
        }
    }
}

%Bib-Datei einbinden
\addbibresource{literatur/literatur.bib}

% Zeilenabstand im Literaturverzeichnis ist Einzeilig
% siehe Leitfaden S. 14
\AtBeginBibliography{\singlespacing}

% -------------------------------------------------------
% Anhangsverzeichnis definieren
% -------------------------------------------------------
\DeclareNewTOC[
  owner    = \jobname,
  listname = {Anhangsverzeichnis},
  tocentryindent   = 1.5em,
  tocentrynumwidth = 2.3em,
]{anh}

\makeatletter
\newcommand*{\useAnhangTOC}{ % \useAnhangTOC wechselt für alle folgenden sections die Aufnahme in das Anhangsverzeichnis
  \renewcommand*{\ext@toc}{anh}%
}
\newcommand{\useMainTOC}{ % \useMainTOC wachselt wieder zurück, sodass alles weitere wieder im Inhaltsverzeichnis erscheint
	\renewcommand*{\ext@toc}{toc}
}
\makeatother

%-----------------------------------
% Silbentrennung (Vgl. Leitfaden S.17) 
%-----------------------------------
\usepackage{hyphsubst}
\HyphSubstIfExists{ngerman-x-latest}{%
\HyphSubstLet{ngerman}{ngerman-x-latest}}{}
\usepackage{microtype}

\hyphenpenalty=100 % reduziert die Trennung
\exhyphenpenalty=2000 % Minimiert die doppelte Trennung

%-----------------------------------
% Pfad fuer Abbildungen
%-----------------------------------
\graphicspath{{./}{./abbildungen/}}

%-----------------------------------
% Weitere Ebene einfügen
%-----------------------------------
\input{skripte/weitereEbene}

%-----------------------------------
% Paket für die Nutzung von Anhängen
%-----------------------------------
\usepackage{appendix}

%-----------------------------------
% Zeilenabstand 1,5-zeilig
%-----------------------------------
\usepackage{setspace}
%\onehalfspacing
\setstretch{1.4} % setstretch statt \onehalfspacing um die Abstände von MS besser nachzuahmen

%-----------------------------------
% Absätze durch eine neue Zeile
%-----------------------------------
\setlength{\parindent}{0mm}
\setlength{\parskip}{0.5em plus 0.3em minus 0.2em}

\sloppy					%Abstände variieren
\pagestyle{headings}

%----------------------------------
% Präfix in das Abbildungs- und Tabellenverzeichnis aufnehmen, statt nur der Nummerierung (siehe Issue #206).
%----------------------------------
\KOMAoption{listof}{entryprefix} % Siehe KOMA-Script Doku v3.28 S.153
\BeforeStartingTOC[lof]{\renewcommand*\autodot{:}} % Für den Doppelpunkt hinter Präfix im Abbildungsverzeichnis
\BeforeStartingTOC[lot]{\renewcommand*\autodot{:}} % Für den Doppelpunkt hinter Präfix im Tabellenverzeichnis
\BeforeStartingTOC[anh]{\renewcommand*\autodot{:}} % Für den Doppelpunkt hinter Präfix im Anhangsverzeichnis
\BeforeStartingTOC[loe]{\renewcommand*\autodot{:}} % Für den Doppelpunkt hinter Präfix im Formelverzeichnis

%-----------------------------------
% Abkürzungsverzeichnis
%-----------------------------------
\usepackage[printonlyused]{acronym}

%-----------------------------------
% Symbolverzeichnis
%-----------------------------------
% Quelle: https://www.namsu.de/Extra/pakete/Listofsymbols.pdf
\usepackage[final]{listofsymbols}

%-----------------------------------
% Glossar
%-----------------------------------
\usepackage{glossaries}
\glstoctrue %Auskommentieren, damit das Glossar nicht im Inhaltsverzeichnis angezeigt wird.
\makenoidxglossaries
\input{abkuerzungen/glossar}

%-----------------------------------
% PDF Meta Daten setzen
%-----------------------------------
\usepackage[hyperfootnotes=true]{hyperref} %Testweise Fußnoten-Links reaktiviert, im Zeifel wieder deaktivieren 
%hyperfootnotes=false deaktiviert die Verlinkung der Fußnote. Ansonsten inkompaibel zum Paket "footmisc"
% Wenn hyperfootnotes=true oder \usepackage{footnotebackref}, dann wird das Komma zwischen zwei an der selben Stelle gesetzten Fußnoten entfernt
% Behebt die falsche Darstellung der Lesezeichen in PDF-Dateien, welche eine Übersetzung besitzen
% siehe Issue 149
\usepackage{footnotebackref} % Dieses Paket setzt links von der Zahl der Fußnote zurück zur Zahl im Text (mit hyperref
\usepackage[all]{hypcap}
\makeatletter
\pdfstringdefDisableCommands{\let\selectlanguage\@gobble}
\makeatother
\ifdefined\myAutorzwei
\hypersetup{
    pdfinfo={
        Title={\myTitel},
        Subject={\myStudiengang},
        Author={\myAutor {, } \myAutorzwei},
        Build=1.1
    }
}
\else
\hypersetup{
    pdfinfo={
        Title={\myTitel},
        Subject={\myStudiengang},
        Author={\myAutor},
        Build=1.1
    }
}
\fi


%-----------------------------------
% PlantUML
%-----------------------------------
%\usepackage{plantuml}

%-----------------------------------
% Umlaute in Code korrekt darstellen
% siehe auch: https://en.wikibooks.org/wiki/LaTeX/Source_Code_Listings
%-----------------------------------
\lstset{literate=
	{á}{{\'a}}1 {é}{{\'e}}1 {í}{{\'i}}1 {ó}{{\'o}}1 {ú}{{\'u}}1
	{Á}{{\'A}}1 {É}{{\'E}}1 {Í}{{\'I}}1 {Ó}{{\'O}}1 {Ú}{{\'U}}1
	{à}{{\`a}}1 {è}{{\`e}}1 {ì}{{\`i}}1 {ò}{{\`o}}1 {ù}{{\`u}}1
	{À}{{\`A}}1 {È}{{\'E}}1 {Ì}{{\`I}}1 {Ò}{{\`O}}1 {Ù}{{\`U}}1
	{ä}{{\"a}}1 {ë}{{\"e}}1 {ï}{{\"i}}1 {ö}{{\"o}}1 {ü}{{\"u}}1
	{Ä}{{\"A}}1 {Ë}{{\"E}}1 {Ï}{{\"I}}1 {Ö}{{\"O}}1 {Ü}{{\"U}}1
	{â}{{\^a}}1 {ê}{{\^e}}1 {î}{{\^i}}1 {ô}{{\^o}}1 {û}{{\^u}}1
	{Â}{{\^A}}1 {Ê}{{\^E}}1 {Î}{{\^I}}1 {Ô}{{\^O}}1 {Û}{{\^U}}1
	{œ}{{\oe}}1 {Œ}{{\OE}}1 {æ}{{\ae}}1 {Æ}{{\AE}}1 {ß}{{\ss}}1
	{ű}{{\H{u}}}1 {Ű}{{\H{U}}}1 {ő}{{\H{o}}}1 {Ő}{{\H{O}}}1
	{ç}{{\c c}}1 {Ç}{{\c C}}1 {ø}{{\o}}1 {å}{{\r a}}1 {Å}{{\r A}}1
	{€}{{\EUR}}1 {£}{{\pounds}}1 {„}{{\glqq{}}}1
}

%-----------------------------------
% Kopfbereich / Header definieren
%-----------------------------------
\pagestyle{fancy}
\fancyhf{}
% Seitenzahl oben, mittig, mit Strichen beidseits
% \fancyhead[C]{-\ \thepage\ -}

% Seitenzahl oben, mittig, entsprechend Leitfaden ohne Striche beidseits
\fancyhead[C]{\thepage}

%\fancyhead[L]{\leftmark}							% kein Footer vorhanden
% Waagerechte Linie unterhalb des Kopfbereiches anzeigen. Laut Leitfaden ist
% diese Linie nicht erforderlich. Ihre Breite kann daher auf 0pt gesetzt werden.
%\renewcommand{\headrulewidth}{0.4pt}
\renewcommand{\headrulewidth}{0pt}

%-----------------------------------
% Damit die hochgestellten Zahlen auch auf die Fußnote verlinkt sind (siehe Issue 169)
%-----------------------------------
\hypersetup{colorlinks=true, breaklinks=true, linkcolor=darkblack, citecolor=darkblack, menucolor=darkblack, urlcolor=darkblack, linktoc=all, bookmarksnumbered=false, pdfpagemode=UseOutlines, pdftoolbar=true}
\urlstyle{same}%gleiche Schriftart für den Link wie für den Text

\usepackage{layout}
%-----------------------------------
% Start the document here:
%-----------------------------------
\begin{document}

\pagenumbering{Roman}								% Seitennumerierung auf römisch umstellen
\newcolumntype{C}{>{\centering\arraybackslash}X}	% Neuer Tabellen-Spalten-Typ:
%Zentriert und umbrechbar

%-----------------------------------
% Textcommands
%-----------------------------------
\input{skripte/textcommands}

%Verzeichnis der Todos

\listoftodos[Anmerkungen]
\newpage

%-----------------------------------
% Titlepage
%-----------------------------------
\begin{titlepage}
	\newgeometry{left=2cm, right=2cm, top=2cm, bottom=2cm}
	\begin{center}
    \includegraphics[width=2.3cm]{abbildungen/fomLogo} \\
    \vspace{.5cm}
		\begin{Large}\textbf{\myHochschulName}\end{Large}\\
    \vspace{.5cm}
		\begin{Large}\langde{Hochschulzentrum}\langen{university location} \myHochschulStandort\end{Large}\\
		\vspace{2cm}
    \begin{large}{\myThesisArt}\end{large}\\
    \vspace{.5cm}

		\langde{im Berufsbegleitenden Studiengang}
		\langen{in the part-time degree program}\\
		\myStudiengang\\
		\vspace{0.5cm}
		\mySemesterZahl. Semester\\
		\vspace{1.7cm}

		\textbf{im Rahmen des Moduls}\\
		\textbf{\myModul}\\
		\vspace{1.8cm}
		\langde{über das Thema}
		\langen{on the subject}\\
    \vspace{0.5cm}
		\large{\textbf{\myTitel}}\\
		\vspace{2cm}
    \langde{von}\langen{by}\\
    \vspace{0.5cm}

\ifdefined\myAutorzwei
\begin{Large}{\myAutor} \& {\myAutorzwei}\end{Large}\\
\else
\begin{Large}{\myAutor}\end{Large}\\
\fi
    
	\end{center}
	\normalsize
	\vfill
    \begin{tabular}{ l l }
        \langde{Betreuer} % für Hausarbeiten
        \langen{Advisor}: & \myBetreuer\\
	\ifdefined\myAutorzwei
		\langde{Matrikelnummern}
        \langen{Matriculation Numbers}: & \myMatrikelNr\ \& \myMatrikelNrzwei\\
	\else
		\langde{Matrikelnummer}
        \langen{Matriculation Number}: & \myMatrikelNr\\
	\fi
        \langde{Abgabedatum}
        \langen{Submission}: & \myAbgabeDatum
    \\
    \end{tabular}
\end{titlepage}


%-----------------------------------
% Vorwort (optional; bei Verwendung beide Zeilen entkommentieren und unter Inhaltsverzeichnis setcounter entsprechend anpassen)
%-----------------------------------
%\input{kapitel/vorwort/vorwort}
%\newpage

%-----------------------------------
% Inhaltsverzeichnis
%-----------------------------------
% Um das Tabellen- und Abbbildungsverzeichnis zu de/aktivieren ganz oben in Documentclass schauen
\setcounter{page}{2}
\addtocontents{toc}{\protect\enlargethispage{-20mm}}% Die Zeile sorgt dafür, dass das Inhaltsverzeichnisseite auf die zweite Seite gestreckt wird und somit schick aussieht. Das sollte eigentlich automatisch funktionieren. Wer rausfindet wie, kann das gern ändern.
\setcounter{tocdepth}{4}
{
\setstretch{1.3} %Zeilenabstand im Inhaltsverzeichnis reduzieren
\tableofcontents
} % Setzt die Zeilenhöhe auf 1,2 statt 1,4 wie es für den Textteil festgelegt ist, um die Zeilenhöhe von 1,5 aus Word möglichst genau wiederzugeben
%\tableofcontents
\newpage

%-----------------------------------
% Abbildungsverzeichnis
%-----------------------------------
\listoffigures
\newpage
%-----------------------------------
% Tabellenverzeichnis
%-----------------------------------
\listoftables
\newpage
%-----------------------------------
% Abkürzungsverzeichnis
%-----------------------------------
% Falls das Abkürzungsverzeichnis nicht im Inhaltsverzeichnis angezeigt werden soll
% dann folgende Zeile auskommentieren.
\addcontentsline{toc}{section}{\abbreHeadingName}
\input{abkuerzungen/acronyms}
\newpage

%-----------------------------------
% Formelverzeichnis
%-----------------------------------
\listofxequations

\vspace{1em}

%-----------------------------------
% Symbolverzeichnis
%-----------------------------------
% In Overleaf führt der Einsatz des Symbolverzeichnisses zu einem Fehler, der aber ignoriert werdne kann
% Falls das Symbolverzeichnis nicht im Inhaltsverzeichnis angezeigt werden soll
% dann folgende Zeile auskommentieren.
%\phantomsection\addcontentsline{toc}{section}{\symheadingname}
\input{skripte/symbolDef}
\makeatletter
\begingroup
  \renewcommand{\symheadingname}{}%
  \let\section\@gobble%
  \vspace{6pt}%
  \listofsymbols
\endgroup
\makeatother
\newpage

%-----------------------------------
% Glossar
%-----------------------------------
\printnoidxglossaries
\newpage

%-----------------------------------
% Sperrvermerk
%-----------------------------------
%\input{kapitel/anhang/sperrvermerk}


%-----------------------------------
% Seitennummerierung auf arabisch und ab 1 beginnend umstellen
%-----------------------------------
\pagenumbering{arabic}
\setcounter{page}{1}

%-----------------------------------
% Kapitel / Inhalte
%-----------------------------------
%TC:endignore - berücksichtigt die folgnden Wörter beim Wordcount in Overleaf
% Die Kapitel werden über folgende Datei eingebunden
\input{skripte/kapitelUebersicht.tex}

%TC:ignore - ignoriert die folgnden Wörter beim Wordcount in Overleaf

%\listoftodos[Anmerkungen]

%-----------------------------------
% Apendix / Anhang
%-----------------------------------
\newpage
\section*{\AppendixName} %Überschrift "Anhang", ohne Nummerierung
\addcontentsline{toc}{section}{\AppendixName} %Den Anhang ohne Nummer zum Inhaltsverzeichnis hinzufügen

% --- Anhangsverzeichnis ---
\ifshowListOfAppendices
    \listofanhs
    \clearpage
\fi

\appendix
\addtocontents{toc}{\protect\setcounter{tocdepth}{3}}
\useAnhangTOC

\begin{appendices}
% Nachfolgende Änderungen erfolgten aufgrund von Issue 163
\makeatletter
\renewcommand\@seccntformat[1]{\csname the#1\endcsname:\quad}
\makeatother
\addtocontents{toc}{\protect\setcounter{tocdepth}{0}} %
	\renewcommand{\thesection}{\AppendixName\ \arabic{section}}
	\renewcommand\thesubsection{\AppendixName\ \arabic{section}.\arabic{subsection}}
	\section{Beispielanhang}\label{Beispielanhang}
Dieser Abschnitt dient nur dazu zu demonstrieren, wie ein Anhang aufgebaut seien kann.
\subsection{Weitere Gliederungsebene}
Auch eine zweite Gliederungsebene ist möglich.
\section{Bilder}
Auch mit Bildern.
Diese tauchen nicht im Abbildungsverzeichnis auf.
\begin{figure}[H]
    \caption[]{Beispielbild}
	\label{fig:Beispielbild}
\begin{center}
    \includegraphics[width=1\textwidth]{verzeichnisStruktur}
\end{center}
    \vspace{-1em}
    %Quelle: Hau bla 2014
    \cite[Quelle:][S. 73]{Wolff2017}
\end{figure}
\end{appendices}
\addtocontents{toc}{\protect\setcounter{tocdepth}{2}}

\useMainTOC

%-----------------------------------
% Literaturverzeichnis
%-----------------------------------
\newpage

% Die folgende Zeile trägt ALLE Werke aus literatur.bib in das
% Literaturverzeichnis ein, egal ob sie zietiert wurden oder nicht.
% Der Befehl ist also nur zum Test der Skripte sinnvoll und muss bei echten
% Arbeiten entfernt werden.
%\nocite{*}

%\addcontentsline{toc}{section}{Literatur}

% Die folgenden beiden Befehle würden ab dem Literaturverzeichnis wieder eine
% römische Seitennummerierung nutzen.
% Das ist nach dem Leitfaden nicht zu tun. Dort steht nur dass 'sämtliche
% Verzeichnisse VOR dem Textteil' römisch zu nummerieren sind. (vgl. S. 3)
%\pagenumbering{Roman} %Zähler wieder römisch ausgeben
%\setcounter{page}{4}  %Zähler manuell hochsetzen

% Ausgabe des Literaturverzeichnisses

% Keine Trennung der Werke im Literaturverzeichnis nach ihrer Art
% (Online/nicht-Online)
%\begin{RaggedRight}
%\printbibliography
%\end{RaggedRight}

% Alternative Darstellung, die laut Leitfaden genutzt werden sollte.
% Dazu die Zeilen auskommentieren und folgenden code verwenden:

% Literaturverzeichnis getrennt nach Nicht-Online-Werken und Online-Werken
% (Internetquellen).
% Die Option nottype=online nimmt alles, was kein Online-Werk ist.
% Die Option heading=bibintoc sorgt dafür, dass das Literaturverzeichnis im
% Inhaltsverzeichnis steht.
% Es ist übrigens auch möglich mehrere type- bzw. nottype-Optionen anzugeben, um
% noch weitere Arten von Zusammenfassungen eines Literaturverzeichnisse zu
% erzeugen.
% Beispiel: [type=book,type=article]
\printbibliography[nottype=online,heading=bibintoc,title={\langde{Literaturverzeichnis}\langen{Bibliography}}]

% neue Seite für Internetquellen-Verzeichnis
\newpage

% Laut Leitfaden 2018, S. 14, Fussnote 44 stehen die Internetquellen NICHT im
% Inhaltsverzeichnis, sondern gehören zum Literaturverzeichnis.
% Die Option heading=bibintoc würde die Internetquelle als eigenen Eintrag im
% Inhaltsverzeicnis anzeigen.
%\printbibliography[type=online,heading=bibintoc,title={\headingNameInternetSources}]
\printbibliography[type=online,heading=subbibliography,title={\headingNameInternetSources}]

\newpage
\section*{KI-Hilfsmittelverzeichnis}

Künstliche Intelligenz wurde im Rahmen der Arbeit zur Quellensuche verwendet.

Folgende Modelle wurden hierfür eingesetzt:

\begin{list}{\bullet}{}
	\item{Claude V4.6 Sonnet}
	\item{GPT-5.4}
\end{list}

\input{kapitel/anhang/erklaerung}
\end{document}

%TC:endignore - benötigt für Wordcount in Overleaf